% !PRINTED 62
\weberpage{62}

We ask first about the possibility of $T$-functions of order zero. If $L(u)$ is such a function and $L'(u)$ its derivative, then twofold differentiation of (6) gives

\[
\frac{d}{du}\,\frac{L'(u+\omega_1)}{L(u+\omega_1)} = \frac{d}{du}\,\frac{L'(u)}{L(u)},
\]

and correspondingly from the second equation (6). Accordingly $d^2\log L(u)/du^2$, being an integral doubly periodic function, is a constant, and consequently $\log L(u)$ is an integral function of the second degree in $u$. We therefore obtain:

4. \gesperrt{A $T$-function of order zero is of the form}

\weq{7}{L(u) = C\,e^{-\pi i (\lambda u^2 + \mu u)},}

where $\lambda$, $\mu$, $C$ are constants.

That every function of this form is a $T$-function of order zero is seen at once.

The product

\weq{8}{T'(u) = L(u)\,T(u)}

is, just like $T(u)$, a $T$-function of order $m$, and vanishes at the same points.

The periodicity factors of $T'$ are obtained from those of $T$ by replacing the exponents $a_1$, $a_2$, $b_1$, $b_2$ by

\weq{9}{\begin{aligned}
a_1' &= a_1 + \lambda\omega_1, & b_1' &= b_1 + \mu\omega_1\\
a_2' &= a_2 + \lambda\omega_2, & b_2' &= b_2 + \mu\omega_2
\end{aligned}}

in which, however, the $b$, $b'$ are determined only up to additive even integers.

5. \gesperrt{By (9), $\mu$ can always be determined, and in only one way, so that $b_1'$, $b_2'$ become real.}

For if the imaginary parts of $b_1'$, $b_2'$ are set equal to zero, one obtains two linear equations for the imaginary and the real part of $\mu$, whose determinant is, by \S\,16, the area of the period parallelogram, and hence different from zero.

The periodicity factors and the zero points of a $T$-function of order $m$ stand in a certain dependence upon one another, which we obtain without difficulty from the theorems of \S\,15. First of all, the order $m$ results if one takes the integral

\[
\frac{1}{2\pi i}\int d\log T(u) = m
\]
%CONT>>
